% SPDX-License-Identifier: CC-BY-4.0
% Copyright (C) 2026 Fabian Januszewski, Abhijit (the Langlandscape team)
% Frame 10 (climax I): the defect formula, statement and status.
% Sources: chapters/s06/s06-m3.tex thm:s06-defect (status as printed there), cor:s06-fixed, cor:s06-allweights (pmc), if:s06-defect (what the talk may claim);
% chapters/s06/s06-m2.tex hyp:s06-H, def:s06-ctam, def:s06-tam; chapters/s06/s06-m8.tex rem:s06-open (i), (iii); chapters/s06/s06-m5.tex prop:s06-levM, rem:s06-levM-status.
% eta_f(M,q) = Kim-Ota's N^- -new congruence ideal (Def. 2.11 of [KO] = eta(T'_k, lambda'_f)); T' = T_B by the Lambda-adic JL (if:s01-targets).
\begin{frame}{The integral defect is a Tamagawa ideal}
  \small
  $D=q$, $N=qM$; $F$ a primitive ordinary $q$-new Hida branch; $\eta_F$, $\eta^B_F$ its congruence ideals on the $\GL_2$ and the $B$ side.
  \begin{block}{Defect formula, under (H0)--(H2), (P), and (H0$^-$) or (Lev)}
    \[\eta_F=\cc_q(F)\cdot\eta^B_F,\qquad (h_F)\;=\;\cc_q(F)\;=\;\tq_q(F)\quad\text{(principal ideals of $I_F$)}.\]
    At each arithmetic point of weight $k\ge3$, trivial character, $q$-new $f=F_P$, \emph{including $k\ge p$}:
    \[\eta_f(N)\;=\;\cc_q(f)\cdot\eta_f(M,q),\qquad \cc_q(f)=\big(\varpi^{\,t_f(q)}\big)\vspace{-2pt}\]
  \end{block}
  $\eta_f(M,q)$: the $q$-new, i.e.\ quaternionic, side. The defect is the old/new congruence ideal, of size $t_f(q)$.\\[1pt]
  {\footnotesize\textbf{Status:} proved modulo cited results under (H0$^-$); under (Lev) for $M=1$ and the cases (F1), (F1$'$) ($N(\bar\rho)=1$; (F1$'$) in our
  sign-fixed reading of Kim--Ota, Thm~3.14). Open: $q\equiv-1\bmod p$, $\bar\rho$ unramified at $q$, outside (Lev) ($N(\bar\rho)>1$, or $\ell^2\mid M$ with $\bar\rho$ unramified at $\ell$).
  Excluded: trivial primes (quaternionic Hecke algebra not Gorenstein \reftag{B\"ockle--Khare--Manning}).}
\end{frame}
